\chapter{Conclusion Générale et Perspectives}

\section{Bilan Technique et Réalisations du Projet}

Ce stage d'ingénieur en 4\textsuperscript{ème} année Génie Informatique, réalisé au sein du Groupe OCP, a permis de concevoir et de développer un prototype fonctionnel et robuste de \textbf{Jumeau Numérique Industriel} appliqué à la supervision et l'analyse d'impact d'une chaîne de production de phosphate.

Les principaux objectifs techniques et méthodologiques atteints au cours de ce projet sont :
\begin{enumerate}
    \item \textbf{Une architecture logicielle pérenne et découplée :} Mise en place rigoureuse de la Clean Architecture et du Domain-Driven Design (DDD) sur le backend .NET 9, garantissant l'indépendance totale du cœur de domaine vis-à-vis des composants d'infrastructure.
    \item \textbf{Une médiation CQRS performante :} Utilisation de MediatR pour isoler les commandes de mutation et les requêtes de lecture.
    \item \textbf{Une analyse d'impact topologique algorithmique :} Développement du service \texttt{ImpactAnalysisService} fondé sur l'algorithme BFS avec détection de cycles, permettant d'identifier automatiquement les cascades de pannes en $\mathcal{O}(|V|+|E|)$.
    \item \textbf{Une supervision temps réel réactive :} Mise en œuvre de WebSockets via ASP.NET Core SignalR couplée à une IHM moderne sous Angular 19 exploitant les \textit{Signals} et la cartographie interactive \texttt{vis-network}.
    \item \textbf{Une haute assurance qualité :} Validation intégrale par une suite de 83 tests unitaires automatisés xUnit/FluentAssertions et 27 fichiers de spécifications frontend Jasmine/Karma.
\end{enumerate}

\section{Apports Personnels et Professionnels}

Sur le plan professionnel, cette immersion au sein du Groupe OCP m'a permis d'appréhender concrètement les contraintes opérationnelles de l'industrie lourde, les exigences de haute disponibilité des chaînes de traitement minier et la nécessité d'une modélisation logicielle rigoureuse.

Sur le plan technique, ce projet a constitué une opportunité d'approfondir la maîtrise de l'écosystème .NET 9 moderne, des patrons de conception d'entreprise (DDD, CQRS, Mediator, Observer), de la programmation réactive sous Angular 19 et de la théorie des graphes appliquée aux processus cyber-physiques.

\section{Perspectives d'Évolution et Feuille de Route}

Ce démonstrateur technologique pose les fondations d'une plateforme de Jumeau Numérique d'entreprise. Les perspectives d'évolution s'articulent selon trois axes :

\begin{itemize}
    \item \textbf{Perspective à Court Terme (Déploiement \& Observabilité) :} Conteneurisation de la solution avec Docker et Docker Compose, et intégration d'outils d'observabilité (OpenTelemetry, Prometheus et tableaux de bord Grafana).
    \item \textbf{Perspective à Moyen Terme (Connectivité Industrielle Standard) :} Remplacement du simulateur de télémétrie par des connecteurs aux protocoles de communication industriels standards (\textbf{OPC-UA}, \textbf{MQTT} ou bus d'événements \textbf{Apache Kafka}) pour ingérer directement les données des automates de l'usine.
    \item \textbf{Perspective à Long Terme (Intelligence Artificielle \& Rendu 3D) :}
    \begin{itemize}
        \item Intégration de modèles de \textit{Machine Learning} pour l'estimation de la durée de vie résiduelle (\textit{Remaining Useful Life --- RUL}) et la maintenance prédictive avancée ;
        \item Développement d'une vue spatiale 3D interactive sous \textbf{Three.js / WebGL} permettant une navigation immersive dans l'atelier de concassage.
    \end{itemize}
\end{itemize}
